% !TEX root = ../main.tex

\chapter{多模态零样本蛋白质功能注释方法MZSGO}

传统的蛋白质功能预测方法将每个GO标签视为独立的分类标识符，模型的预测能力被限制在训练集中出现过的标签范围内，无法对新出现的功能术语进行有效推理。针对这一根本性局限，本章提出多模态零样本蛋白质功能注释方法MZSGO（Multimodal Zero-Shot GO annotation），将蛋白质功能预测重新定义为开放语义空间中的"蛋白质-GO"匹配问题。MZSGO融合蛋白质语言模型提取的序列进化特征与大语言模型编码的文本语义特征，通过自适应门控融合机制实现多模态信息的动态整合，使模型既能在有监督场景下准确分类，也能通过文本语义相似性推断未见过的功能标签。

\section{问题定义}

蛋白质功能预测任务的目标是：给定一个蛋白质$P$和一组GO功能标签$\mathcal{G} = \{g_1, g_2, \ldots, g_C\}$，预测该蛋白质与每个功能标签之间的关联概率$p(y=1 | P, g_c)$，其中$y=1$表示蛋白质$P$具有功能$g_c$。

与传统方法不同，本文不将其建模为固定标签集的多标签分类问题，而是重新定义为"蛋白质-GO"匹配问题。具体而言，每个蛋白质的特征表示与每个GO标签的文本语义嵌入配对，由模型计算二者的匹配分数。这种范式设计意味着：对于训练集中未出现过的新GO标签，只需提供其文本定义即可生成语义嵌入并参与预测，从而天然支持零样本推理。

考虑到GO的三个本体（BP、MF、CC）在标签规模、语义粒度和数据分布上存在显著差异，本文对每个本体独立训练一个模型。这一策略基于如下实验观察：联合训练所有本体会导致模型难以学习各本体特有的特征模式，独立训练能够获得更好的性能。

\section{数据集构建}

为严格评估模型在有监督和零样本两个场景下的性能，本文构建了基于CAFA5挑战和Swiss-Prot数据库\cite{boutet2007uniprotkb}的综合基准数据集。

\subsection{有监督训练集构建}

训练语料源自CAFA5数据集，采用ProtGO\cite{wang2025protgo}使用的数据源和训练/测试划分策略，并限制数据条目为Swiss-Prot数据库中收录的蛋白质。这一子集由人工审核的蛋白质组成，确保了注释质量的高可靠性，同时为与需要结构信息的基线方法进行公平比较提供了基础。功能注释采用2023年1月1日发布的Gene Ontology版本。

与ProtGO不同的是，本文实施了本体特异性的数据过滤策略。对于每个本体类别，仅使用具有该特定本体注释的蛋白质进行训练，而非直接使用完整训练集。这一决策基于如下观察：完整数据集中缺乏某一本体注释的蛋白质会对基线模型的性能产生不利影响。尽管MZSGO在完整数据集上也不会出现显著性能下降，但经过裁剪的训练集有助于实现最优性能并确保与最先进基线方法的严格公平比较。

\subsection{零样本测试集构建}

零样本测试集旨在评估模型对从未见过的蛋白质和新出现标签的注释能力。其设计遵循两个核心原则：

（1）\textbf{蛋白质新颖性（序列层面）}：选择不在训练集中的Swiss-Prot蛋白质。为防止序列相似性导致的信息泄露，使用Diamond工具进行严格的同源性过滤，移除与训练集中任何序列的序列一致性超过30\%的测试蛋白质。

（2）\textbf{标签新颖性（语义层面）}：实施时间演进评估策略，引入2025年10月10日更新的GO本体版本。将"零样本标签"定义为存在于2025版本但在2023版本中不存在的GO术语。因此，最终的零样本测试集由具有这些新出现功能注释的非同源蛋白质组成。

\subsection{数据集统计}

数据集的统计信息如表~\ref{tab:mzsgo_dataset_stats}所示，样本按三个GO子本体分组。由于生物通路的复杂性，BP类别包含最多的零样本蛋白质数量。

\begin{table}[h]
\centering
\caption{MZSGO数据集统计}
\label{tab:mzsgo_dataset_stats}
\begin{tabular}{lccc}
\toprule
数据集 & MF & CC & BP \\
\midrule
训练集 & 50,874 & 51,743 & 51,618 \\
测试集 & 5,584 & 5,764 & 5,712 \\
零样本测试集 & 555 & 17 & 3,687 \\
\bottomrule
\end{tabular}
\end{table}

表~\ref{tab:mzsgo_label_stats}详细展示了标签分布情况。在测试阶段，标签被分为以下三类：共有标签（Shared Labels），即同时出现在训练集和测试集中的GO术语；未见标签（Unseen Labels），即存在于测试集但不在训练数据中的术语；零样本标签（Zero-shot Labels），特指因本体更新而新出现的术语。

\begin{table}[h]
\centering
\caption{标签分布统计}
\label{tab:mzsgo_label_stats}
\begin{tabular}{lccc}
\toprule
标签类别 & MF & CC & BP \\
\midrule
标签总数 & 7,007 & 2,803 & 20,639 \\
未见标签 & 191 & 26 & 328 \\
共有标签 & 2,983 & 1,581 & 11,691 \\
仅训练标签 & 3,833 & 1,196 & 8,620 \\
零样本标签 & 35 & 6 & 34 \\
\bottomrule
\end{tabular}
\end{table}

这种严格的区分确保了零样本测试集上的评估指标真正反映模型的推理能力，而非记忆效应。

\section{MZSGO方法设计}

\subsection{总体框架}

MZSGO的总体架构如图~\ref{fig:mzsgo_architecture}所示。系统由四个紧密联结的组件构成：（1）多模态特征编码，提取蛋白质序列、域描述和GO标签定义的嵌入表示；（2）特征投影与非对称特征丢弃，将异构表示映射至共享语义空间并增强训练鲁棒性；（3）自适应门控融合，动态整合多模态特征；（4）分类预测，输出蛋白质与GO功能的匹配概率。

\begin{figure}[h]
\centering
\includegraphics[width=1\textwidth]{figures/model.pdf}
\caption{MZSGO框架总体架构。流程包含四个集成阶段：（1）多模态特征投影；（2）训练阶段的蛋白质模态丢弃，虚线表示以概率$p=0.15$随机丢弃蛋白质特征，实线表示保留的模态；（3）门控融合模块实现自适应特征整合；（4）二分类头进行GO功能预测。}
\label{fig:mzsgo_architecture}
\end{figure}

\subsection{多模态特征编码}

MZSGO利用两种预训练模型提取三个模态的特征表示：

（1）\textbf{蛋白质序列表示}：使用ESM2-650M\cite{lin2023evolutionary}蛋白质语言模型对蛋白质序列进行编码，提取最后一层隐藏状态，通过平均池化获得蛋白质级别的全局进化特征表示$\mathbf{E}_{seq} \in \mathbb{R}^{1280}$。

（2）\textbf{GO标签语义表示}：使用Qwen3-Embedding-4B\cite{zhang2025qwen3}对每个GO术语的名称和定义文本进行编码。输入文本格式为"名称: 定义"（如"protein kinase activity: Catalysis of the phosphorylation of an amino acid residue..."），通过平均池化获得标签级别的语义嵌入$\mathbf{E}_{label} \in \mathbb{R}^{2560}$。

（3）\textbf{蛋白质域语义表示}：使用同一Qwen3-Embedding-4B模型对蛋白质各个域的InterPro\cite{blum2025interpro}描述文本进行编码，采用两级平均池化（token级和域级）获得蛋白质级别的域特征表示$\mathbf{E}_{dom} \in \mathbb{R}^{2560}$。

以上三种特征分别捕获了蛋白质的进化信号、功能标签的语义内涵和蛋白质域的结构功能信息，为后续的多模态融合提供了互补的信息来源。

\subsection{特征投影与模态特异性丢弃}

由于三种模态的原始嵌入具有不同的维度和分布特性（ESM嵌入为1280维，文本嵌入为2560维），需要将它们映射至统一的隐空间。为此，设计三个独立的多层感知机（MLP）投影模块，每个模块由线性变换、层归一化\cite{ba2016layer}、GELU激活\cite{hendrycks2016gaussian}和Dropout组成：
\begin{equation}
    \mathbf{h}_m = \text{Dropout}(\text{GELU}(\text{LN}(\mathbf{W}_m \mathbf{E}_m + \mathbf{b}_m)))
\end{equation}
其中$m \in \{seq, dom, label\}$表示模态类型，$\mathbf{h}_m \in \mathbb{R}^{d_{model}}$为投影后的特征向量，$d_{model} = 512$。

\textbf{模态特异性丢弃策略}。为防止模型过度依赖单一蛋白质模态，并模拟真实场景中域信息缺失的情况，本文设计了一种非对称的特征丢弃机制。在训练阶段，定义二值掩码$\mathbf{M} = [m_{seq}, m_{dom}]$，其中每个蛋白质模态以概率$p = 0.15$被独立丢弃：
\begin{equation}
    m_{seq} = \mathbb{I}[r_{seq} > p], \quad m_{dom} = \mathbb{I}[r_{dom} > p]
\end{equation}
其中$r_{seq}, r_{dom} \sim \text{Uniform}(0, 1)$。掩码后的特征为$\hat{\mathbf{h}}_{seq} = \mathbf{h}_{seq} \cdot m_{seq}$和$\hat{\mathbf{h}}_{dom} = \mathbf{h}_{dom} \cdot m_{dom}$。

关键设计在于其"非对称性"：GO标签的文本嵌入$\mathbf{h}_{label}$始终保留，不参与丢弃。这迫使模型学会在部分蛋白质信息缺失的情况下，仍能通过GO标签的文本语义进行有效预测，从而为零样本泛化奠定基础。

为确保至少有一个蛋白质模态可用，设计了强制保留机制：
\begin{equation}
    \text{若 } m_{seq} = m_{dom} = 0: \quad m_k \leftarrow 1, \quad k \sim \text{Uniform}\{seq, dom\}
\end{equation}
推理阶段不施加任何丢弃，模型可利用完整的多模态信息进行预测。

\subsection{自适应门控融合机制}

为有效整合多模态特征，本文设计了门控融合模块（Gated Fusion Module），学习为每个信息来源分配自适应的重要性权重。首先，拼接三个模态的特征构建全局上下文向量，输入门控网络生成注意力分数：
\begin{equation}
    \mathbf{S} = [\hat{\mathbf{h}}_{dom} \parallel \hat{\mathbf{h}}_{seq} \parallel \mathbf{h}_{label}]
\end{equation}
\begin{equation}
    \boldsymbol{\alpha} = \text{Softmax}(\mathbf{W}_{gate} \cdot \text{ReLU}(\text{LN}(\mathbf{W}_{ctx}\mathbf{S})))
\end{equation}
其中$\parallel$表示拼接操作，$\boldsymbol{\alpha} \in \mathbb{R}^{3}$为归一化的门控权重，$\mathbf{W}_{ctx} \in \mathbb{R}^{d_{model} \times 3d_{model}}$和$\mathbf{W}_{gate} \in \mathbb{R}^{3 \times d_{model}}$为可学习参数。

同时，每个模态的特征通过各自的变换层$\phi(\cdot)$进行非线性映射。最终的融合表示通过加权求和获得：
\begin{equation}
    \mathbf{H}_{fused} = \sum_{k=1}^{3} \alpha_k \cdot \phi(\mathbf{h}_k)
\end{equation}
其中$k$索引三个模态（序列、域、标签），$\phi(\cdot)$包含线性变换、层归一化、ReLU激活和Dropout。

该机制的核心优势在于：门控权重是\textbf{逐样本动态计算}的，即每个蛋白质-GO对的融合比例不同。例如，当蛋白质的域信息丰富且与功能高度相关时，模型倾向于分配更高的权重给域特征；当域信息稀疏时，模型则更依赖序列特征和标签文本语义。

\subsection{优化目标}

融合表示$\mathbf{H}_{fused}$通过分类MLP生成标量逻辑值，再经sigmoid激活函数获得预测概率：
\begin{equation}
    \hat{y} = \sigma(\mathbf{W}_{out}\mathbf{H}_{fused} + b_{out})
\end{equation}

为应对不同GO子本体在类别不平衡程度上的差异，本文对不同本体采用不同的损失函数策略：

对于BP和MF子本体，使用标准的二分类交叉熵损失：
\begin{equation}
    \mathcal{L}_{\text{BCE}} = -\frac{1}{N} \sum_{i=1}^{N} \left[ y_i \log(\hat{y}_i) + (1 - y_i) \log(1 - \hat{y}_i) \right]
\end{equation}

对于CC子本体，采用Focal Loss以聚焦于难样本的学习：
\begin{equation}
    \mathcal{L}_{\text{Focal}} = -\frac{1}{N} \sum_{i=1}^{N} \left[ \alpha (1 - \hat{y}_i)^\gamma y_i \log(\hat{y}_i) + (1 - \alpha) \hat{y}_i^\gamma (1 - y_i) \log(1 - \hat{y}_i) \right]
\end{equation}
其中$\gamma = 2$为聚焦参数。

\section{实验设置}

\subsection{实现细节}

MZSGO基于PyTorch\cite{paszke2019pytorch}框架实现。为对齐异构模态，将1280维ESM嵌入和2560维文本/域特征投影至统一的隐空间（$d_{model} = 512$）。这些特征经拼接后输入门控融合网络——一个两层MLP——计算自适应模态权重，随后通过变换层进行非线性映射。融合表示最终被压缩至256维用于分类。

训练在单块NVIDIA A6000 GPU（48GB显存）上进行，使用Adam优化器\cite{kingma2014adam}，学习率为$5 \times 10^{-4}$，批次大小为16。为缓解过拟合，施加0.5的神经元级Dropout率。考虑到各子本体的复杂度差异，BP训练30个epoch，MF和CC各训练20个epoch。

\subsection{基线方法}

本文选择十种先进的基线方法进行对比，涵盖三大类别：

（1）\textbf{传统方法}：Naive（基于标签频率预测）、Diamond\cite{diamond}（基于序列同源性传递）。

（2）\textbf{深度学习方法}：DeepGOCNN\cite{kulmanov2020deepgoplus}（卷积神经网络）、DeepGOZero\cite{kulmanov2022deepgozero}（本体公理约束）、Struct2GO\cite{jiao2023struct2go}（图神经网络+结构信息）、PO2GO\cite{li2024partial}（偏序嵌入）、DeepGO-SE\cite{kulmanov2024protein}（神经符号推理）。

（3）\textbf{多模态/LLM方法}：ProtGO\cite{wang2025protgo}（四模态融合）、DPFunc\cite{wang2025dpfunc}（域信息增强）、ProtNote\cite{char2025protnote}（LLM功能描述）。

\section{实验结果与分析}

\subsection{标准测试集上的性能对比}

表~\ref{tab:mzsgo_comparison}汇总了MZSGO与所有基线方法在三个本体上的性能对比结果。

\begin{table*}[htbp]
\centering
\caption{MZSGO与最先进方法在测试集上的性能对比。最优结果加粗，次优结果加下划线。}
\label{tab:mzsgo_comparison}
\resizebox{\textwidth}{!}{%
\begin{tabular}{llccccc}
\toprule
\textbf{本体} & \textbf{模型} & \textbf{Fmax}$\uparrow$ & \textbf{AUPR}$\uparrow$ & \textbf{Unseen AUPR}$\uparrow$ & \textbf{Seen AUPR}$\uparrow$ & \textbf{H}$\uparrow$ \\
\midrule
\multirow{11}{*}{\textbf{BP}} 
 & Naive & 0.2991 & 0.2073 & 0.0025 & 0.2076 & 0.0049 \\
 & Diamond & 0.4703 & 0.2913 & 0.0025 & 0.2934 & 0.0050 \\
 & DeepGOCNN & 0.3964 & 0.3676 & 0.0025 & 0.3684 & 0.0050 \\
 & DeepGOZero & 0.4888 & 0.4535 & 0.0025 & 0.4558 & 0.0050 \\
 & Struct2GO & 0.4022 & 0.3744 & 0.0025 & 0.3747 & 0.0050 \\
 & PO2GO & 0.2975 & 0.2303 & 0.0025 & 0.2307 & 0.0049 \\
 & DeepGO-SE & \underline{0.5006} & 0.4373 & 0.0025 & 0.4404 & 0.0050 \\
 & ProtGO & 0.4636 & 0.4377 & 0.0025 & 0.4410 & 0.0050 \\
 & DPFunc & 0.4984 & \underline{0.4902} & 0.0069 & \underline{0.4936} & 0.0136 \\
 & ProtNote & 0.3555 & 0.2993 & \underline{0.0303} & 0.2997 & \underline{0.0550} \\
 & \textbf{MZSGO} & \textbf{0.5045} & \textbf{0.5008} & \textbf{0.2393} & \textbf{0.5021} & \textbf{0.3241} \\
\midrule
\multirow{11}{*}{\textbf{MF}} 
 & Naive & 0.5496 & 0.3318 & 0.0032 & 0.3327 & 0.0063 \\
 & Diamond & 0.6727 & 0.3330 & 0.0032 & 0.3366 & 0.0063 \\
 & DeepGOCNN & 0.6603 & 0.5225 & 0.0032 & 0.5242 & 0.0064 \\
 & DeepGOZero & 0.7441 & 0.6333 & 0.0032 & 0.6380 & 0.0064 \\
 & Struct2GO & 0.6564 & 0.5940 & 0.0032 & 0.5962 & 0.0064 \\
 & PO2GO & 0.5496 & 0.4136 & 0.0032 & 0.4148 & 0.0064 \\
 & DeepGO-SE & 0.7481 & 0.6332 & 0.0032 & 0.6376 & 0.0064 \\
 & ProtGO & 0.7262 & \underline{0.7399} & 0.0032 & \underline{0.7453} & 0.0064 \\
 & DPFunc & \underline{0.7495} & \textbf{0.7416} & 0.0051 & \textbf{0.7480} & 0.0101 \\
 & ProtNote & 0.6140 & 0.5283 & \underline{0.0490} & 0.5301 & \underline{0.0897} \\
 & \textbf{MZSGO} & \textbf{0.7611} & 0.7337 & \textbf{0.4806} & 0.7380 & \textbf{0.5821} \\
\midrule
\multirow{11}{*}{\textbf{CC}} 
 & Naive & 0.6143 & 0.0793 & 0.0207 & 0.5270 & 0.0398 \\
 & Diamond & 0.6331 & 0.3076 & 0.0207 & 0.3093 & 0.0388 \\
 & DeepGOCNN & 0.6918 & 0.5284 & 0.0207 & 0.5292 & 0.0398 \\
 & DeepGOZero & 0.7075 & 0.6598 & 0.0207 & 0.6612 & 0.0401 \\
 & Struct2GO & 0.6989 & 0.7380 & 0.0207 & 0.7382 & 0.0403 \\
 & PO2GO & 0.6143 & 0.5828 & 0.0207 & 0.5830 & 0.0400 \\
 & DeepGO-SE & 0.7379 & 0.6991 & 0.0207 & 0.7003 & 0.0402 \\
 & ProtGO & 0.7209 & \underline{0.7717} & 0.0207 & \underline{0.7727} & 0.0403 \\
 & DPFunc & \underline{0.7386} & 0.7709 & 0.0274 & 0.7722 & 0.0529 \\
 & ProtNote & 0.6636 & 0.6605 & \underline{0.2530} & 0.6607 & \underline{0.3659} \\
 & \textbf{MZSGO} & \textbf{0.7470} & \textbf{0.7726} & \textbf{0.5862} & \textbf{0.7733} & \textbf{0.6669} \\
\bottomrule
\end{tabular}%
}
\end{table*}

\subsubsection{有监督性能分析}

在常规分类指标（Fmax和AUPR）方面，MZSGO在三个本体上均取得了最优或极具竞争力的结果，凸显了多模态融合方法的有效性。

在最具挑战性的BP本体中，MZSGO以0.5045的Fmax超过了域感知方法DPFunc（0.4984）和神经符号推理方法DeepGO-SE（0.5006）。在MF和CC本体上，MZSGO分别以0.7611和0.7470的Fmax保持领先。值得注意的是，MZSGO在CC本体上的AUPR达到0.7726，刷新了该本体的最优记录。

与仅依赖单一模态的方法相比，MZSGO通过自适应门控融合机制有效整合了蛋白质序列特征（ESM）与域描述语义，使模型能够综合利用互补的语义线索，产生更强的特征表示。

\subsubsection{零样本性能分析}

MZSGO的核心优势在于其对未见标签的推理能力。表~\ref{tab:mzsgo_comparison}中的Unseen AUPR和调和均值H直观地反映了这一点。

传统比对方法和标准深度学习模型在零样本场景中表现出极为有限的泛化能力。这主要是因为这些模型将标签视为独立的one-hot标识符，在训练过程中，优化目标自然地抑制了未见标签对应的权重，迫使模型对训练集中未出现的类别预测接近零的概率。值得注意的是，尽管DeepGOZero理论上能够通过逻辑约束实现零样本预测，但它在缺乏显式GO公理的罕见类别上预测失效，Unseen AUPR接近0。

ProtNote利用LLM展示了一定的零样本能力（如BP上的0.0303），但MZSGO大幅超越了它。具体而言，MZSGO在BP本体上实现了0.2393的Unseen AUPR——相较ProtNote提升了约\textbf{8倍}。在MF和CC本体上同样观察到一致的趋势，MZSGO分别实现了0.4806和0.5862的Unseen AUPR。

综合平衡已见和未见标签性能的调和均值H进一步凸显了模型的全面优势。MZSGO在BP、MF、CC上分别获得0.3241、0.5821和0.6669的H值，大幅超越亚军ProtNote的0.0550、0.0897和0.3659。这一结果证实MZSGO不仅仅是记忆训练标签，而是学到了蛋白质特征与功能描述之间的泛化性语义对齐。

\subsection{纯零样本测试集上的评估}

为进一步严格验证泛化能力，本文在专门构建的纯零样本测试集上将MZSGO与零样本能力最强的基线方法ProtNote进行直接对比。结果如表~\ref{tab:mzsgo_zeroshot}所示。

\begin{table}[htbp]
\centering
\caption{纯零样本测试集上的性能对比。该评估集通过时间演进方式构建，评估模型对2025版本中新出现但训练截止日期前不存在的GO功能标签的预测能力。}
\label{tab:mzsgo_zeroshot}
\begin{tabular}{llcccc}
\toprule
\textbf{本体} & \textbf{模型} & \textbf{Fmax} & \textbf{AUPR} & \textbf{Precision} & \textbf{Recall} \\
\midrule
\multirow{2}{*}{BP} 
 & ProtNote & 0.3763 & 0.2705 & 0.2580 & \textbf{0.6951} \\
 & \textbf{MZSGO} & \textbf{0.4881} & \textbf{0.4546} & \textbf{0.4475} & 0.5369 \\
\midrule
\multirow{2}{*}{MF} 
 & ProtNote & 0.3599 & 0.3393 & 0.4545 & 0.2979 \\
 & \textbf{MZSGO} & \textbf{0.5982} & \textbf{0.5435} & \textbf{0.6261} & \textbf{0.5727} \\
\midrule
\multirow{2}{*}{CC} 
 & ProtNote & 0.2983 & 0.1652 & 0.1853 & 0.7647 \\
 & \textbf{MZSGO} & \textbf{0.6610} & \textbf{0.6291} & \textbf{0.5521} & \textbf{0.8235} \\
\bottomrule
\end{tabular}
\end{table}

MZSGO在所有三个本体的Fmax和AUPR上均一致超越ProtNote。对精确率和召回率的精细分析揭示了两种方法的行为差异和改进方向：

（1）\textbf{BP本体的平衡精确率-召回率权衡}：ProtNote展现了较高的召回率（0.6951）但精确率低下（0.2580），表明其倾向于过度预测功能标签。相比之下，MZSGO实现了更均衡的权衡，将精确率提升至0.4475的同时保持了0.5369的竞争性召回率，从而获得了显著更高的Fmax（0.4881 vs. 0.3763）。

（2）\textbf{MF和CC本体的全面优势}：在MF本体上，MZSGO获得0.5982的Fmax（ProtNote为0.3599），精确率（0.6261 vs. 0.4545）和召回率（0.5727 vs. 0.2979）双双领先。在CC本体上，性能差距最为显著：ProtNote的精确率仅为0.1853，可能是由于广泛文本匹配导致的"功能幻觉"；而MZSGO达到0.5521的精确率，同时召回率也优于ProtNote（0.8235 vs. 0.7647），实现了Fmax从0.2983到0.6610的翻倍提升。

这一对比结果表明，MZSGO通过将预测锚定在域特异性的语义空间中，有效避免了其他基于文本的模型中常见的噪声和假阳性问题，为蛋白质新功能注释提供了高度可靠的预测能力。

\subsection{消融实验}

为评估各组件的贡献——特别是多模态特征、融合模块和优化策略，本文在CC本体上进行了系统的消融实验。消融变体分为五组：（1）单模态编码器；（2）特征表示变化；（3）架构变化；（4）编码器替换；（5）损失函数变体。结果如表~\ref{tab:mzsgo_ablation}所示。

\begin{table}[htbp]
  \centering
  \caption{CC本体上的消融实验结果。报告指标包括总体Fmax和AUPR，以及零样本相关指标。H为已见与未见标签性能的调和均值。}
  \label{tab:mzsgo_ablation}
  \resizebox{\linewidth}{!}{
    \begin{tabular}{lccccc}
    \toprule
    \textbf{方法} & \textbf{Fmax} & \textbf{AUPR} & \textbf{Unseen AUPR} & \textbf{Unseen Fmax} & \textbf{H} \\
    \midrule
    \multicolumn{6}{l}{\textit{模态消融}} \\
    仅ESM特征 & 0.7319 & 0.7344 & 0.5203 & 0.0793 & 0.6094 \\
    仅域特征 & 0.7311 & 0.7052 & 0.5187 & 0.2322 & 0.6012 \\
    \midrule
    \multicolumn{6}{l}{\textit{特征与架构消融}} \\
    域信息One-hot编码 & 0.7203 & 0.6731 & 0.5047 & 0.2000 & 0.5889 \\
    使用PO2Vec标签嵌入 & 0.7406 & 0.7196 & 0.3193 & 0.1818 & 0.4423 \\
    去除特征丢弃 & \textbf{0.7544} & 0.6857 & 0.4898 & 0.2500 & 0.5714 \\
    简单拼接替代门控融合 & 0.7381 & 0.7009 & 0.5245 & 0.2469 & 0.6000 \\
    \midrule
    \multicolumn{6}{l}{\textit{编码器变体}} \\
    BioGPT & 0.7419 & 0.7198 & 0.4258 & 0.1935 & 0.5351 \\
    BiomedBERT & 0.7406 & 0.7218 & 0.5647 & 0.2000 & 0.6337 \\
    Qwen3-0.6B & 0.7406 & 0.7282 & 0.5469 & 0.1379 & 0.6247 \\
    ESM-150M & 0.7385 & 0.7310 & 0.3864 & 0.1333 & 0.5056 \\
    \midrule
    \multicolumn{6}{l}{\textit{损失函数变体}} \\
    BCE损失 & 0.7501 & 0.7143 & 0.5824 & 0.2581 & 0.6421 \\
    加权BCE损失 & 0.7199 & 0.7450 & 0.5127 & 0.4444 & 0.6095 \\
    \midrule
    \textbf{MZSGO} & 0.7470 & \textbf{0.7726} & \textbf{0.5862} & \textbf{0.5665} & \textbf{0.6669} \\
    \bottomrule
    \end{tabular}%
  }
\end{table}

从消融结果中可以得出以下关键结论：

（1）\textbf{域语义信息是零样本能力的关键}。移除域语义信息造成最严重的性能退化：仅使用ESM特征时，Unseen Fmax从0.5665骤降至0.0793，表明序列特征本身难以泛化到未见标签。仅使用域特征时Unseen Fmax为0.2322，远优于仅用ESM；而将域信息退化为One-hot编码后Unseen Fmax降至0.2000，证实零样本迁移依赖于域描述的语义内容而非简单标识符。

（2）\textbf{门控融合机制不可或缺}。将门控融合替换为简单拼接使H从0.6669降至0.6000，说明朴素拼接无法有效整合异构模态。使用PO2Vec标签嵌入替代文本嵌入更是导致H降至0.4423，这是因为基于图拓扑的PO2Vec嵌入与序列和文本模态的语义对齐度较低。

（3）\textbf{特征丢弃策略增强泛化能力}。去除特征丢弃后，已见标签准确度略有提升（Fmax达0.7544），但零样本泛化显著下降（H从0.6669降至0.5714），验证了其作为正则化手段的重要性。

（4）\textbf{编码器选择的合理性}。各替代编码器的性能均不及当前组合，证实MZSGO选择的编码器（ESM2-650M + Qwen3-Embedding-4B）最能捕获跨模态的互补结构信息。

（5）\textbf{Focal Loss在CC本体上的优势}。标准BCE获得较高的Fmax（0.7501）但总体AUPR偏低（0.7143）；加权BCE虽改善了Unseen Fmax（0.4444）但损害了总体AUPR和H值。Focal Loss（MZSGO）通过有效平衡困难样本的学习，实现了最优的综合权衡。

\subsection{参数敏感性分析}

为评估模型的鲁棒性，本文考察了两个关键超参数的影响：Dropout率和隐藏向量维度。采用控制变量法，每次仅变动一个参数，在CC和MF本体上观察AUPR的波动。

\subsubsection{Dropout率的影响}

Dropout率应用于特征提取后的神经元级丢弃，取值范围为0.3至0.7。模型在Dropout率为0.5时表现最佳，CC本体AUPR达0.7726，MF本体达0.7377。较低的Dropout率（如0.3）导致过拟合和泛化能力下降；较高的Dropout率（如0.7）因过度信息丢失使性能退化（CC降至0.7610，MF降至0.7228）。因此，0.5被选为最优配置。

\subsubsection{隐藏向量维度的影响}

隐藏维度从128变化至1024。随着维度从128扩展至512，两个本体的性能持续提升。CC本体在512维时达到峰值（AUPR为0.7726），此后随维度增大而退化。MF本体在1024维时仍有提升（AUPR达0.8133），但CC本体在该维度下显著下降至0.7371。为构建跨GO类别通用的模型架构而无需针对特定任务进行调整，最终将隐藏维度设定为512，在表征能力和计算成本之间取得了良好的平衡。

\section{本章小结}

本章提出了多模态零样本蛋白质功能注释方法MZSGO，其核心创新在于三个方面：（1）将蛋白质功能预测从封闭标签集分类重新定义为开放语义空间中的蛋白质-GO匹配问题，通过大语言模型构建域和标签的统一语义空间，从根本上赋予模型零样本推理能力；（2）设计了自适应门控融合模块，实现逐样本动态的多模态特征整合，克服了传统简单拼接方式的局限；（3）引入非对称特征丢弃策略，在训练阶段增强模型对模态缺失的鲁棒性。

实验结果表明，MZSGO在三个GO本体的有监督指标上均达到最优或极具竞争力的水平，在零样本场景下则以压倒性优势超越所有现有方法。其中，在最具挑战性的CC本体零样本场景下，MZSGO的调和均值H达到0.6669，是亚军方法ProtNote（0.3659）的近两倍。消融实验验证了域语义信息、门控融合机制和特征丢弃策略各自的有效性。MZSGO为蛋白质功能注释，特别是新功能标签的预测，提供了可靠且有效的计算工具。
